%% classes -- tree classes: two seeds, observation pins, finite-bound pinning
%% for every variance, regular and finite transfer, the top-freeness gadget
%% for {bot,f}, and arities zero and one.

\section{Tree Classes}\label{sec:tc}\label{sec:classes}

The preceding sections decide existential entailment over arbitrary trees.
This section transfers the decision to regular and to finite trees, for every
arity $n\ge0$ and every variance $c$, and it reduces the signature $\BF$ with
$f$ covariant to $\BFT$. Regular trees are handled by the least closures of
\cref{sec:lf}, whose slices are regular for regular seeds by
\cref{lem:sm-regular-closure}; a failed extension is caused by at most two
seeds (\cref{lem:tc-two-seeds}), so $\Ereg=\Eall$ (\cref{thm:tc-regular}).
Finite trees need a different argument. Least closures describe trees with as
few leaves as possible, and these are infinite in general: the empty system has the empty least model, which
assigns $\Omega$ to every variable. Instead we express observations by finite
ground bounds (\cref{lem:tc-pins}) and show that adding finite ground bounds
preserves finite satisfiability, for every variance (\cref{thm:tc-pinning}).
This gives the dispatch \cref{eq:tc-finite} of \cref{thm:tc-finite}. Proofs
not given here are in \cref{app:tc}.

Throughout, trees are over $\BFT$ and ordered by $\le_c$ unless stated
otherwise, $\varphi$ and $\psi$ are as in \cref{def:bg-existential}, we use
the notation of \cref{sec:bg,sec:lf} together with $f(t^n)=f(t,\dots,t)$ and
$\nu_\pm$ for $\nu_0$ and $\nu_1$, and a model $T$ \emph{decodes} to the
solution $\rho$ with $S_\rho=T$ (\cref{lem:lf-clauses}).

\subsection{Extensions and Two Seeds}\label{sec:tc-seeds}

An assignment $\rho$ of $X$ is \emph{extensible} if it extends to a solution
of $\psi$ over arbitrary trees.

\begin{lemma}[Extension criterion]\label{lem:tc-extension}
An assignment $\rho$ of $X$ is extensible iff $H_\psi(S_\rho)$ is consistent
and $H_\psi(S_\rho)|_X=S_\rho$. In that case $H_\psi(S_\rho)$ decodes to an
extension of $\rho$.
\end{lemma}

This restates \cref{lem:lf-closure}(5) for assignments.

\begin{proof}
By \cref{lem:lf-closure}(5) and \cref{lem:lf-clauses}, applied to the set
$S_\rho$, which is extension closed and has no exclusion pair, $\rho$ is
extensible exactly when $H_\psi(S_\rho)$ is consistent and has visible part
$S_\rho$; then that closure decodes to an extension.
\end{proof}

\begin{definition}[Observations]\label{def:tc-observation}
A \emph{visible observation} is a visible atom $a=(x^\pm,w)$ together with one
of the words \emph{asserted} or \emph{denied}. An asserted observation
\emph{holds} in an assignment $\rho$ of $X$ if $a\in S_\rho$, a denied one if
$a\notin S_\rho$. The variable $x$ and the word $w$ are the variable and the
word of the observation.
\end{definition}

\begin{lemma}[Two seeds]\label{lem:tc-two-seeds}
Let $\rho$ be an assignment of $X$ that is not extensible. There are a set
$A\subseteq S_\rho$ of visible atoms and a set $C$ of visible atoms outside
$S_\rho$, with $\len A+\len C\le2$, such that no assignment $\rho'$ of $X$
with $A\subseteq S_{\rho'}$ and $C\cap S_{\rho'}=\varnothing$ is extensible.
In particular, at most two visible observations hold in $\rho$ and hold
together in no extensible assignment.
\end{lemma}

\begin{proof}
Every atom $b$ of $H_\psi(S_\rho)$ has a seed in $\Fact_\psi\cup S_\rho$, so
$b\in H_\psi(A_b)$ for a set $A_b\subseteq S_\rho$ with at most one element. By
\cref{lem:tc-extension}, $H_\psi(S_\rho)$ is not consistent or has a visible
atom outside $S_\rho$. In the first case, \cref{lem:lf-closure}(3) gives
seeds $x,y\in\Fact_\psi\cup S_\rho$ such that $\Reach_\psi(x)$ contains a
forbidden root, or $\Reach_\psi(x)$ and $\Reach_\psi(y)$ contain an exclusion pair;
let $A$ consist of those of $x,y$ that lie in $S_\rho$, and $C=\varnothing$.
In the second case, let $b\in H_\psi(S_\rho)|_X$ with $b\notin S_\rho$, and put
$A=A_b$ and $C=\{b\}$. In both cases $\len A+\len C\le2$. Let $\rho'$ be an
assignment of $X$ with $A\subseteq S_{\rho'}$ and $C\cap
S_{\rho'}=\varnothing$. Since $H_\psi$ is monotone,
$H_\psi(S_{\rho'})\supseteq H_\psi(A)$, so $H_\psi(S_{\rho'})$ contains the
same forbidden root or the same two opposite atoms, or the visible atom
$b\notin S_{\rho'}$. By \cref{lem:tc-extension}, $\rho'$ is not extensible.
The observations asserting the atoms of $A$ and denying those of $C$ prove the
last claim.
\end{proof}

The lemma uses no satisfiability hypothesis. If $\psi$ is unsatisfiable, then
$\mu_\psi$ is inconsistent by \cref{lem:lf-closure}(2), and
$A=C=\varnothing$ will do. The two seeds of $A$ may belong to the same
variable.

\subsection{Observation Pins}\label{sec:tc-pins}

For a word $w$ and trees $b,g$, let $\emb_b(w,g)$ be the tree with a pin path
along $w$ that carries $g$ at its end and $b$ at every child off the pin path:
$\emb_b(\eps,g)=g$, and $\emb_b(iw,g)=f(t_j)_j$ with $t_i=\emb_b(w,g)$ and
$t_j=b$ for $j\ne i$.

\begin{lemma}[Observation pins]\label{lem:tc-pins}
For every covariant tree $t$ and every word $w$,
\begin{equation}\label{eq:tc-pins}
\begin{aligned}
 w\in L^0_t&\iff\emb_\Bot(w,\Top)\le t,&
 w\notin L^0_t&\iff t\le\emb_\Top(w,f(\Top^n)),\\
 w\in L^1_t&\iff t\le\emb_\Top(w,\Bot),&
 w\notin L^1_t&\iff\emb_\Bot(w,f(\Bot^n))\le t,
\end{aligned}
\end{equation}
and each of these four trees has at most $1+n(\len w+1)$ positions.
Consequently, for every variance $c$, every visible observation with variable
$x$ and word $w$ is equivalent, for all assignments $\rho$ of $X$, to a single
inequality $g\le_c\rho(x)$ or $\rho(x)\le_cg$, where $g$ is a finite ground
tree with at most $1+n(\len w+1)$ positions.
\end{lemma}

For $n\ge1$ the bound is attained by the two trees with a constructor at the
end of the pin path; for $n=0$ the only word is $\eps$, and the four trees are
$\Top$, $f()$, $\Bot$ and $f()$.

\subsection{Finite-Bound Pinning}\label{sec:tc-pinning}

\paragraph{Intersection of leaf cones}
Over arbitrary and regular trees, extensions can be intersected; over finite
trees they cannot (item (3) below), which is why the transfer to finite trees
pins observations by finite ground bounds instead. For trees $t$ and $u$, let
$t\leafmeet u$ be the tree with the leaf pair
$(L^0_t\cap L^0_u,\ L^1_t\cap L^1_u)$, which exists by
\cref{thm:lf-representation}(1); its signed leaf cones contain exactly the
observations common to $t$ and $u$. For assignments, $\leafmeet$ acts variablewise.

\begin{proposition}[Closure under intersection]\label{prop:ch-meet}
\leavevmode
\begin{enumerate}[leftmargin=*]
\item The intersection of a nonempty family of models of a system $\chi$ is a
  model of $\chi$, and the intersection of a nonempty subfamily of the
  projected family $\mathcal F_\psi$ (\cref{eq:ch-models}) belongs to
  $\mathcal F_\psi$. In particular, the solutions of $\chi$ and the assignments of $X$ that extend
  to solutions of $\psi$ are closed under $\leafmeet$.
\item The operation $\leafmeet$ is commutative, associative, idempotent and
  monotone for $\le_c$, and
  $f(t_i)_i\leafmeet f(u_i)_i=f(t_i\leafmeet u_i)_i$. Moreover,
  $t\leafmeet u$ is the median of $t$, $u$ and $\Omega$ in the lattice of
  trees ordered by $\le_c$.
\item If $t$ and $u$ are regular, then so is $t\leafmeet u$. For $n\ge1$,
  $\Bot\leafmeet\Top=\Omega$ is infinite, so the finite trees are not
  closed under $\leafmeet$.
\end{enumerate}
\end{proposition}

The proof is in \cref{app:tc}. By (1), the model $H_\psi(S)$ of
\cref{lem:lf-closure}(5) is the intersection of all models $T$ of $\psi$ with
$T|_X=S$, the least extension of $S$. By (1) and commutativity, a relation
that is the projection of the solutions of a system over arbitrary trees, or
over regular trees, and contains two pairs $(a,b)$ and $(b,a)$ also contains
the diagonal pair $(a\leafmeet b,a\leafmeet b)$. So no such relation is
nonempty, symmetric and irreflexive; in particular, over arbitrary or regular
trees no system with hidden variables defines the disequality of two trees.
The decision procedure does not use \cref{prop:ch-meet}.

\paragraph{Ground bounds}
A \emph{ground bound} on a variable $x$ is an inequality $g\le_cx$, a
\emph{lower} bound, or $x\le_cg$, an \emph{upper} bound, where $g$ is a finite
ground tree. The following theorem extends
\nssecite{Theorem ``Finite countermodel transfer''} from the pin paths of a
single word to arbitrary finite bounds.

\begin{theorem}[Finite-bound pinning]\label{thm:tc-pinning}
Let $\chi$ be a system over a set $V$ of variables, with any variance, that
has a finite solution, and let $B$ be a finite set of ground bounds on
variables of $V$. If $\chi\cup B$ has a solution, then it has a finite
solution.
\end{theorem}

The proof reduces the theorem to covariant systems by the signed system
$\chi^\pm$ of \cref{tab:bg-signed}, and then symmetrizes a covariant solution
by a median. A ground bound $b$ is translated into a pair $b^\pm$ of
covariant ground bounds on signed variables,
\begin{equation}\label{eq:tc-signed-bounds}
 (g\le_cx)^\pm=\{\nu_+g\le x^+,\ x^-\le\nu_-g\},\qquad
 (x\le_cg)^\pm=\{x^+\le\nu_+g,\ \nu_-g\le x^-\},
\end{equation}
and $B^\pm$ is the union of these pairs. For an assignment $\rho$ of $V$,
$\rho^\pm$ is the assignment $v^\pm\mapsto\nu_\pm\rho(v)$ of
\cref{lem:bg-signed}. The \emph{sign involution} $\iota$ maps an assignment
$\sigma$ of $V^\pm$ to the assignment
\[
 (\iota\sigma)(v^+)=\sigma(v^-)^\partial,\qquad
 (\iota\sigma)(v^-)=\sigma(v^+)^\partial .
\]
The literals of \cref{tab:bg-signed} that contain $\Bot$ or $\Top$ belong to
$\chi^\pm$; the proof needs them (\cref{rem:tc-extreme}).

\begin{lemma}[Signed bounds]\label{lem:tc-signed}
Let $\chi$ be a system over $V$ and $B$ a finite set of ground bounds on $V$.
\begin{enumerate}[leftmargin=*]
\item An assignment $\rho$ of $V$ solves $\chi\cup B$ iff $\rho^\pm$ solves
  $\chi^\pm\cup B^\pm$.
\item $\iota$ is an involution that maps the solutions of $\chi^\pm\cup B^\pm$
  to solutions, and $\iota\rho^\pm=\rho^\pm$ for every assignment $\rho$ of
  $V$.
\item If $\iota\sigma=\sigma$, then $\sigma=\rho^\pm$ for the assignment
  $\rho(v)=\nu_+\sigma(v^+)$; and $\rho\mapsto\rho^\pm$ and its inverse
  preserve finiteness.
\end{enumerate}
\end{lemma}

For covariant trees $t_1,t_2,t_3$, the \emph{median} $\tcmed(t_1,t_2,t_3)$ is
the covariant tree whose cones consist of the words that lie in the
corresponding cones of at least two of the three trees: for both signs,
\[
 L^\pm_{\tcmed(t_1,t_2,t_3)}=\{w:w\in L^\pm_{t_j}\text{ for at least two }j\}.
\]
It exists by \cref{thm:lf-representation}(1): the sets on the right are
extension closed, and they are disjoint, since a word in two of the cones
$L^0_{t_j}$ and in two of the cones $L^1_{t_j}$ lies in both cones of one
$t_j$. The median of assignments is taken variablewise. It is the lattice
median $(t_1\wedge t_2)\vee(t_1\wedge t_3)\vee(t_2\wedge t_3)$ of the
covariant order, but we use only the following properties.

\begin{lemma}[Median]\label{lem:tc-median}
\leavevmode
\begin{enumerate}[leftmargin=*]
\item The median is symmetric and monotone in each argument, and
  $\tcmed(t,t,u)=t$.
\item $\tcmed(f(a^1_i)_i,f(a^2_i)_i,f(a^3_i)_i)=f(\tcmed(a^1_i,a^2_i,a^3_i))_i$
  and $\tcmed(t_1,t_2,t_3)^\partial=\tcmed(t_1^\partial,t_2^\partial,t_3^\partial)$.
  In particular $\iota$ commutes with medians of assignments of signed
  variables.
\item The median of three finite trees is finite.
\item If three assignments solve a covariant system together with covariant
  ground bounds, so does their median.
\end{enumerate}
\end{lemma}

The covariant case of \cref{thm:tc-pinning} is proved directly from the
finite-satisfiability criterion of \citet{Zhao26}
(\cref{thm:bg-nsse-sat}), in \cref{app:tc-covariant}. By that
criterion, a system has a finite solution iff it has neither a label clash nor
a cycle clash. The proof encodes each ground bound by fresh variables, one for
each position of the bound; since a path through fresh variables strictly
increases the length of the position they stand for, every cycle clash of the
encoded system lies among the original variables, where the finite solution
of $\chi$ excludes it, and the solution of $\chi\cup B$ excludes label
clashes.

\begin{lemma}[Covariant finite-bound pinning]\label{lem:tc-pinning-covariant}
\Cref{thm:tc-pinning} holds when $f$ is covariant.
\end{lemma}

\begin{proof}[Proof of \cref{thm:tc-pinning}]
If $n=0$, every tree is finite. Let $n\ge1$, let $\delta$ be a finite solution
of $\chi$, and let $\rho$ be a solution of $\chi\cup B$. By
\cref{lem:tc-signed}(1,3), $\delta^\pm$ is a finite solution of $\chi^\pm$,
and $\rho^\pm$ is a solution of $\chi^\pm\cup B^\pm$. By
\cref{lem:tc-pinning-covariant}, applied to the covariant system $\chi^\pm$
and the bounds $B^\pm$, the system $\chi^\pm\cup B^\pm$ has a finite solution
$\alpha$. By \cref{lem:tc-signed}(2), $\iota\alpha$ is a solution of
$\chi^\pm\cup B^\pm$, and it is finite because duality keeps the domain. Consider the
median
\[
 \tcmed(\alpha,\iota\alpha,\delta^\pm).
\]
All three arguments solve $\chi^\pm$, so $\tcmed(\alpha,\iota\alpha,\delta^\pm)$
solves $\chi^\pm$ by \cref{lem:tc-median}(4), and it is finite by
\cref{lem:tc-median}(3). The assignment $\delta^\pm$ need not satisfy the
bounds; still $\tcmed(\alpha,\iota\alpha,\delta^\pm)$ satisfies $B^\pm$: for a
lower bound $G\le p$ in $B^\pm$, both $\alpha(p)$ and $(\iota\alpha)(p)$ lie
above $G$, so
$G=\tcmed(G,G,\delta^\pm(p))\le\tcmed(\alpha,\iota\alpha,\delta^\pm)(p)$ by
\cref{lem:tc-median}(1), and upper bounds are symmetric. By
\cref{lem:tc-median}(2) and \cref{lem:tc-signed}(2),
\[
 \iota\,\tcmed(\alpha,\iota\alpha,\delta^\pm)
 =\tcmed(\iota\alpha,\iota\iota\alpha,\iota\delta^\pm)
 =\tcmed(\iota\alpha,\alpha,\delta^\pm)=\tcmed(\alpha,\iota\alpha,\delta^\pm) .
\]
By \cref{lem:tc-signed}(3), $\tcmed(\alpha,\iota\alpha,\delta^\pm)=\tau^\pm$ for
a finite assignment $\tau$ of $V$, which solves $\chi\cup B$ by \cref{lem:tc-signed}(1).
\end{proof}

\begin{remark}[The extreme literals are needed]\label{rem:tc-extreme}
The signed literals of $u=\Bot$ and $\Top\le u$ in \cref{tab:bg-signed} cannot
be dropped from $\chi^\pm$, although they contribute no rules on leaf cones,
only facts (\cref{tab:lf-clauses}). Let $n\ge1$, $\chi=\{\Top\le x\}$,
$B=\varnothing$, and $\delta(x)=\Top$. Without its extreme literals the signed
system is empty, and $\alpha(x^+)=\Bot$, $\alpha(x^-)=\Top$ is a finite
solution with $\iota\alpha=\alpha$. Then
$\tcmed(\alpha,\iota\alpha,\delta^\pm)=\alpha$, which decodes to $x=\Bot$ and
violates $\chi$. The third argument of the median must also be a finite
solution of the unpinned system: $\Omega$ in every coordinate is fixed by
$\iota$ and solves every signed system without extreme literals, but the
median with it is infinite in general.
\end{remark}

\subsection{Regular and Finite Trees}\label{sec:tc-transfer}

\begin{theorem}[Regular trees]\label{thm:tc-regular}
For every arity $n\ge0$, every variance $c$, and all systems $\varphi$ over
$X$ and $\psi$ over $X\cup Z$:
\begin{enumerate}[leftmargin=*]
\item a system is satisfiable iff it has a regular solution;
\item every regular extensible assignment of $X$ has a regular extension;
\item if some solution of $\varphi$ is not extensible, then some regular
  solution of $\varphi$ is not extensible.
\end{enumerate}
Consequently $\Ereg(\varphi,\psi)\iff\Eall(\varphi,\psi)$.
\end{theorem}

\begin{proof}
(1) This restates \cref{thm:bg-nsse}(1) and \cref{lem:bg-signed-sat}.

(2) Let $\rho$ be regular and extensible. By \cref{lem:tc-extension},
$H_\psi(S_\rho)$ decodes to an extension of $\rho$. The cones of $\rho$ are
regular by \cref{thm:lf-representation}(4), hence so are those of
$H_\psi(S_\rho)$ by \cref{lem:sm-regular-closure}, and the extension is
regular by \cref{thm:lf-representation}(4).

(3) Let $\rho$ be a solution of $\varphi$ that is not extensible, and take $A$
and $C$ from \cref{lem:tc-two-seeds}. The model $S_\rho$ of $\varphi$ contains
$A$, so by \cref{lem:lf-closure}(2) the set $H_\varphi(A)$ is the least model
of $\varphi$ that contains $A$, and $H_\varphi(A)\subseteq S_\rho$. It decodes
to a solution $\rho_0$ of $\varphi$, which is regular by
\cref{lem:sm-regular-closure} and \cref{thm:lf-representation}(4), since $A$
is finite. Its atom set contains
$A$ and, being contained in $S_\rho$, is disjoint from $C$. By
\cref{lem:tc-two-seeds}, $\rho_0$ is not extensible.

For the consequence, if $\Eall$ holds, every regular solution of $\varphi$ is
extensible and has a regular extension by (2); if $\Eall$ fails, (3) gives a
regular solution of $\varphi$ without any extension.
\end{proof}

\begin{theorem}[Finite trees]\label{thm:tc-finite}
For every arity $n\ge0$, every variance $c$, and all systems $\varphi$ over
$X$ and $\psi$ over $X\cup Z$,
\begin{equation}\label{eq:tc-finite}
 \Efin(\varphi,\psi)\iff
 \neg\Sat_{\mathrm{fin}}(\varphi)\ \lor\
 \bigl(\Sat_{\mathrm{fin}}(\psi)\land\Eall(\varphi,\psi)\bigr).
\end{equation}
More precisely:
\begin{enumerate}[leftmargin=*]
\item if $\varphi$ has a finite solution and some solution of $\varphi$ is
  not extensible, then some finite solution of $\varphi$ is not extensible;
\item if $\psi$ has a finite solution, then every finite extensible assignment
  of $X$ has a finite extension.
\end{enumerate}
\end{theorem}

\begin{proof}
(1) Let $\rho$ be a solution of $\varphi$ that is not extensible. By
\cref{lem:tc-two-seeds}, at most two visible observations hold in $\rho$ and
hold together in no extensible assignment. By \cref{lem:tc-pins}, each of them
is equivalent to a ground bound on a visible variable; let $B$ be the set of
these bounds. Then $\rho$ solves $\varphi\cup B$, and $\varphi$ has a finite
solution, so \cref{thm:tc-pinning} gives a finite solution $\rho_1$ of
$\varphi\cup B$. The observations hold in $\rho_1$, so $\rho_1$ is not
extensible.

(2) Let $\rho$ be a finite assignment of $X$ with an extension $\rho'$ that
solves $\psi$. Let $B$ consist of the bounds $\rho(x)\le_cx$ and
$x\le_c\rho(x)$ for $x\in X$. Then $\rho'$ solves $\psi\cup B$, and $\psi$ has
a finite solution, so \cref{thm:tc-pinning} gives a finite solution $\tau$ of
$\psi\cup B$. By antisymmetry of $\le_c$, $\tau$ agrees with $\rho$ on $X$, so
it is a finite extension of $\rho$.

For \cref{eq:tc-finite}: if $\varphi$ has no finite solution, both sides
hold. If $\varphi$ has a finite solution and $\psi$ has none, that solution
has no finite extension, and both sides fail. If both have finite solutions,
(1) and (2) give $\Efin(\varphi,\psi)\iff\Eall(\varphi,\psi)$.
\end{proof}

The bounds in the proof of (1) concern at most two visible variables, and each
has at most $1+n(\len w+1)$ positions, where $w$ is the word of its
observation. The proofs of \cref{thm:tc-regular,thm:tc-finite} use no
satisfiability hypothesis and no condition on the variance.

\begin{corollary}[Decisions for regular and finite trees]\label{cor:tc-decidable}
For fixed $n$ and $c$, if existential entailment over $\BFT$ is decidable over
arbitrary trees, then it is decidable over regular trees and over finite
trees: $\Ereg=\Eall$, and $\Efin$ is decided by \cref{eq:tc-finite}, whose
two finite-satisfiability tests are decidable by \cref{thm:bg-nsse}(1).
\end{corollary}

\begin{remark}[Both finite-satisfiability tests are needed]\label{rem:tc-examples}
Let $n\ge1$ and let $c$ be any variance. For $X=\varnothing$, the empty
antecedent and $\psi=\{z\le f(z^n),\ f(z^n)\le z\}$, the entailment holds over
arbitrary and regular trees, with the unique witness $z=\Omega$, but fails
over finite trees, since $\psi$ has no finite solution. For
$\varphi=\{x\le f(x^n),\ f(x^n)\le x\}$ and $\psi=\{x=\Bot\}$, the entailment
fails over arbitrary and regular trees, at $x=\Omega$, and holds vacuously
over finite trees. Least models are not finite countermodels in general, and
by \cref{prop:ch-meet}(3) the finite trees are not closed under intersection
of leaf cones; this is why \cref{thm:tc-finite} pins observations instead.
\end{remark}

\subsection{Trees without Top}\label{sec:tc-gadget}

In this subsection $f$ is covariant, so $\le_c$ is the covariant order $\le$.
A tree is \emph{$\Top$-free} if no position is labeled $\Top$; the trees over
$\BF$ are the $\Top$-free trees (\cref{sec:bg-trees}), and a $\Top$-free tree
is finite or regular in the same sense over $\BF$ and over $\BFT$. We write
$\models^{\BF}_{\mathcal D}$ and $\models^{\BFT}_{\mathcal D}$ for existential
entailment over the two signatures, with $f$ covariant; over $\BF$ the systems
have no literal $\Top\le u$ and all values are $\Top$-free
(\cref{def:bg-existential}). We use the trees $T_0=\Bot$ and
$T_{k+1}=f(T_k^n)$ of \cref{sec:bg-constraints}, now for every $n\ge0$: the
positions of $T_k$ are the words of length at most $k$, labeled $f$ below
length $k$ and $\Bot$ at length $k$, and for $n=0$ and $k\ge1$ the tree $T_k$
is $f()$. For a finite tree $t$, $h(t)$ is the maximal length of a position of
$t$. The \emph{gadget} on a variable $v$, with a fresh
variable $y_v$, is the pair of literals
\[
 \mathsf{TF}(v,y_v)=\{v\le y_v,\ y_v\le f(y_v^n)\}.
\]

\begin{lemma}[Top-freeness gadget]\label{lem:tc-gadget}
Let $f$ be covariant.
\begin{enumerate}[leftmargin=*]
\item If $y\le f(y^n)$, then $y$ is $\Top$-free, and so is every tree $t\le y$.
\item For every $\Top$-free tree $t$, the regular tree $y=\Omega$ satisfies
  $t\le y\le f(y^n)$. If $t$ is finite, then so does the finite tree
  $y=T_{h(t)+1}$.
\end{enumerate}
\end{lemma}

\begin{proof}
(1) For covariant trees $L^0_t$ is the set of words with a prefix labeled
$\Top$ in $t$, since $\nu_+t=t$; so $t$ is $\Top$-free iff $L^0_t=\varnothing$.
By \cref{thm:lf-representation}(2,3), $y\le f(y^n)$ gives
\[
 L^0_y\subseteq L^0_{f(y^n)}=\bigcup_{i\in\Sigma}i\,L^0_y=\Sigma L^0_y .
\]
If $L^0_y$ were nonempty, a shortest word in it would have the form $iv$ with
$v\in L^0_y$, contradicting minimality. So $L^0_y=\varnothing$; for $n=0$ this
holds because $\Sigma L^0_y=\varnothing$. If $t\le y$, then
$L^0_t\subseteq L^0_y=\varnothing$ by \cref{thm:lf-representation}(2). The
argument concerns finite words only and applies to infinite trees $y$.

(2) Every position of $\Omega$ is labeled $f$, so $t(w)\le\Omega(w)$ at every
common position $w$ since $t$ is $\Top$-free; and $\Omega=f(\Omega^n)$ has a
single subtree. If $t$ is finite, then $\Kpos(t)$ consists of words of length
less than $h(t)+1$, so $\Kpos(t)\subseteq\Kpos(T_{h(t)+1})$ and $t\le
T_{h(t)+1}$ by \cref{lem:bg-bf-order}. By the same lemma,
$T_{h(t)+1}\le T_{h(t)+2}=f(T_{h(t)+1}^n)$.
\end{proof}

\begin{theorem}[Reduction of $\BF$ to $\BFT$]\label{thm:tc-gadget}
Let $f$ be covariant, and let $\varphi$ over $X$ and $\psi$ over $X\cup Z$ be
systems over $\BF$. Take fresh variables $y_v$ for all $v\in X\cup Z$, put
$Y_X=\{y_x:x\in X\}$ and $Y_Z=\{y_z:z\in Z\}$, and let
\[
 \varphi^\natural=\varphi\cup\bigcup_{x\in X}\mathsf{TF}(x,y_x),\qquad
 \psi^\natural=\psi\cup\bigcup_{z\in Z}\mathsf{TF}(z,y_z),
\]
with visible variables $X\cup Y_X$ and hidden variables $Z\cup Y_Z$. For every
tree class $\mathcal D$,
\[
 \varphi\models^{\BF}_{\mathcal D}\Exists\psi
 \iff
 \varphi^\natural\models^{\BFT}_{\mathcal D}\exists\bar z\,\bar y\,\psi^\natural ,
\]
where $\bar z\,\bar y$ lists $Z\cup Y_Z$.
\end{theorem}

The gadget is placed on every visible variable $x\in X$, including those that
do not occur in $\varphi$, and on every hidden variable $z\in Z$; the visible
variables of $\psi$ are covered by the gadgets of $\varphi^\natural$.

\begin{proof}
For a $\Top$-free tree $t$ in $\mathcal D$ let $y(t)=T_{h(t)+1}$ if
$\mathcal D=\Trees_{\mathrm{fin}}$ and $y(t)=\Omega$ otherwise. By
\cref{lem:tc-gadget}(2), $y(t)$ is a $\Top$-free tree in $\mathcal D$ with
$t\le y(t)\le f(y(t)^n)$. A $\Top$-free assignment satisfies a literal over
$\BF$ in the sense of $\BF$ iff it satisfies it in the sense of $\BFT$, since
both signatures use the same trees, constructor and order on $\Top$-free
trees (\cref{sec:bg-trees}).

Assume the left side, and let $\rho$ be a solution of $\varphi^\natural$ with
values in $\mathcal D$. By \cref{lem:tc-gadget}(1), the values $\rho(y_x)$ and
$\rho(x)$, $x\in X$, are $\Top$-free. So $\rho|_X$ is a solution of $\varphi$
over $\BF$ with values in $\mathcal D$, and it has an extension $\sigma$ to
$Z$, with $\Top$-free values in $\mathcal D$, that solves $\psi$. Put
$\sigma(y_z)=y(\sigma(z))$ for $z\in Z$. Then $\rho$ together with $\sigma$
solves $\psi^\natural$, which does not constrain the variables $y_x$.

Assume the right side, and let $\rho$ be a solution of $\varphi$ over $\BF$
with values in $\mathcal D$. Extend it by $\rho(y_x)=y(\rho(x))$ for $x\in X$;
this gives a solution of $\varphi^\natural$ over $\BFT$ with values in
$\mathcal D$. It has an extension $\sigma$ to $Z\cup Y_Z$, with values in
$\mathcal D$, that solves $\psi^\natural$. By \cref{lem:tc-gadget}(1), every
$\sigma(y_z)$ and every $\sigma(z)$ is $\Top$-free, so $\sigma|_Z$ extends
$\rho$ to a solution of $\psi$ over $\BF$ with values in $\mathcal D$.
\end{proof}

\begin{corollary}[The signature $\BF$]\label{cor:tc-bf}
Let $f$ be covariant, of any arity $n\ge0$. In each tree class, existential
entailment over $\BF$ reduces to existential entailment over $\BFT$ in the same
tree class, at the same arity and with $f$ covariant. The reduction doubles the
numbers of visible and of hidden variables, so the size parameter $s$ becomes
$2s$, and it adds two literals per variable. Consequently, existential
entailment over $\BF$ with $f$ covariant is decidable over arbitrary, regular
and finite trees whenever the covariant problem over $\BFT$ is, and
$\Ereg=\Eall$ holds over $\BF$ as well.
\end{corollary}

The last claim follows from \cref{thm:tc-gadget,thm:tc-regular}. The finite
witness $T_{h(t)+1}$ can be large, but it is used only in the proof; the
reduction introduces variables and literals, not ground trees. Both
requirements on the gadget are necessary.

\begin{remark}[Every visible variable, and covariance]\label{rem:tc-gadget-needed}
The gadget cannot be omitted on a visible variable that does not occur in
$\varphi$: with $X=\{x\}$, the entailment
$\mathrm{true}\models^{\BF}_{\mathcal D}\exists\bar z\ x\le f(z_i)_i$ holds in
every tree class (take the $z_i$ to be the children of $x$, or $\Bot$ if
$x=\Bot$), whereas without a gadget on $x$ its translation fails at
$x=\Top$. The gadget also needs covariance. If $c(j)=1$ for some $j$, let
$y=f(t_i)_i$ with $t_i=\Top$ if $c(i)=1$ and $t_i=\Bot$ if $c(i)=0$. Then
$y\le_cf(y^n)$, since the covariant arguments require $\Bot\le_cy$ and the
contravariant ones $y\le_c\Top$, but $y$ has a label $\Top$; for the arrow,
$y=\Top\to\Bot$ satisfies $y\le_c y\to y$. We do not treat $\BF$ with a
contravariant argument. That problem is not obtained from $\BFT$ by omitting
the literals $\Top\le u$: for every variance, the entailment above holds over
$\BF$ but fails over $\BFT$ at $x=\Top$.
\end{remark}

\subsection{Arities Zero and One}\label{sec:tc-small}

\begin{remark}[Arity zero]\label{rem:tc-nullary}
Let $n=0$. By \cref{rem:bg-nullary}, all three tree classes coincide at arity
zero, so $\Efin=\Ereg=\Eall$, and existential entailment is decidable by
finite enumeration; each of the $3^s$ assignments is checked in time linear in
$\len\varphi+\len\psi$. The route of \cref{sec:lf,sec:ch} also applies, with
no nonempty words, but the enumeration is simpler. Copy literals cannot be
eliminated at arity zero (\cref{lem:bg-sandwich} needs $n\ge1$), which is why
the proof of \cref{thm:tc-pinning} treats $n=0$ separately.
\end{remark}

\begin{remark}[Arity one]\label{rem:tc-unary}
Let $n=1$. Every tree is $f^k(\Bot)$ or $f^k(\Top)$ for some $k\ge0$, or
$\Omega$, so every tree is regular and \cref{thm:tc-regular} is immediate.
Finite and arbitrary trees still differ: the system
$\{x\le f(x),\ f(x)\le x\}$ has the unique solution $\Omega$. The variance
$c(0)=1$ is allowed, and then mixed polarity already occurs: for
$\{x\le f(x)\}$ there is no parity assignment $\pi$ as in
\cref{rem:ch-parity}, since the condition $\pi(a)=\pi(u)\oplus c(i)$ for a
constructor literal with root $u$ and child $a$ at argument $i$ would require
$\pi(x)=\pi(x)\oplus1$. The spine automaton of \cref{sec:nm} runs over
$\Sigma=\{0\}$, so no state has two distinct safe direction letters. Hence its
branching case does not occur, and every infinite direction word is $0^\omega$.
The general argument applies without change.
\end{remark}
